\subsection{Reduced algebras and etale algebras}

DEFINITION 2. --- Let A be a commutative ring; then A is said to be reduced if every nilpotent element (I, p.~98) of A is zero.

If A is a field, or an integral domain, or a product of reduced rings, it is a reduced ring. For a commutative ring A to be reduced it is necessary and sufficient that \(a^{2} \neq 0\) for every \(a \neq 0\) in A: for from this we obtain by induction on n, \(a^{2^{n}} \neq 0\) , hence \(a^{n} \neq 0\) for any \(a \neq 0\) in A.

An algebra is said to be reduced if its underlying ring is reduced.

PROPOSITION 5. --- Let A be a commutative algebra of finite degree over K. For A to be reduced it is necessary and sufficient that there exist extensions \(L_{1}, \ldots, L_{n}\) of finite degree over K such that A is K-isomorphic to \(L_{1} \times \cdots \times L_{n}\) .

The stated condition is clearly sufficient.

Conversely, assume that A is reduced; arguing by induction on the degree of A we see that it suffices to prove that if A is not a field, there exist two non-zero algebras \(A_{1}\) and \(A_{2}\) such that A is isomorphic to \(A_{1} \times A_{2}\) , or also that there exists in A an idempotent different from 0 and 1.

Suppose from now on that A is reduced and is not a field. Among the ideals of A different from 0 and A let a be an ideal whose dimension as vector K-space is minimal. For any \(x \neq 0\) in a we have \(x^{2} \neq 0\) , because A is reduced, whence \(a^{2} \neq \{0\}\) . We have \(a^{2} \subset a\) and by the minimality of a we find that \(a^{2} = a\) . By Lemma 4 there exists an idempotent e such that a = Ae, and we have \(e \neq 0\) , \(e \neq 1\) because a is distinct from 0 and A.

THEOREM 4. --- Let A be a commutative algebra of finite degree over K. Then the following assertions are equivalent :

\begin{enumerate}
\def\labelenumi{\alph{enumi})}
\item
  The algebra \(\mathbf{A}\) is etale.
\item
  For every extension L of K, the ring L ⊗K A is reduced.
\item
  There exists a perfect extension field \(\mathbf{P}\) of \(\mathbf{K}\) such that the ring \(\mathbf{P} \otimes_{\mathbf{K}} \mathbf{A}\) is reduced.
\end{enumerate}

* d) There exist separable algebraic extensions \(\mathbf{L}_1, \ldots, \mathbf{L}_n\) of \(\mathbf{K}\) such that \(\mathbf{A}\) is isomorphic to \(\mathbf{L}_1 \times \cdots \times \mathbf{L}_n\) . *

In particular, every etale algebra is reduced.

\begin{enumerate}
\def\labelenumi{\Alph{enumi})}
\tightlist
\item
  Let us first prove the equivalence of \(a\) , \(b\) and \(c\) .
\end{enumerate}

Suppose that A is etale and let L be an extension of K. Let \(\Omega\) be an algebraically closed extension field of L (V, p.~23, Th. 2). Then \(L \otimes_{K} A\) is isomorphic to a subring of \(\Omega \otimes_{K} A\) and the latter is isomorphic to a ring \(\Omega^{n}\) by Prop. 2 (V, p.~30). Therefore the ring \(L \otimes_{K} A\) is reduced.

We have thus shown that \(a)\) implies \(b)\) , and \(c)\) is a particular case of \(b)\) . Assume now that \(c)\) holds. For the K-algebra A to be etale it is necessary and sufficient that the P-algebra \(\mathbf{A}_{(P)}\) should be etale (V, p.~32, Cor. 2). The algebra A is therefore etale by the following lemma:

Lemma 5. --- Let B be a reduced algebra of finite degree over a perfect field P; then B is etale.

By Prop. 5 there exist extensions \(L_{1}, \ldots, L_{n}\) of P such that B is isomorphic to the algebra \(L_{1} \times \cdots \times L_{n}\) . Since a finite product of etale algebras is etale (V, p.~31, Cor.) it is enough to examine the case where B is an extension of P. By Th. 3 (V, p.~33) it is enough to show that dx = 0 in \(\Omega_{\mathrm{P}}(\mathrm{B})\) for all \(x \in B\) .

Let \(x \in B\) ; since \(B\) is of finite degree over \(K\) , \(x\) is algebraic over \(K\) (V, p.~18, Prop. 2). Let \(f\) be the minimal polynomial of \(x\) and let \(f'\) be the derivative of \(f\) . The polynomial \(f\) is non-constant. Suppose that \(f' = 0\) ; by \(V\) , p.~9, Cor. the field \(P\) is of characteristic \(p \neq 0\) and we have \(f \in P[X^p]\) ; since \(P\) is perfect, we have \(P[X^p] = P[X]^p\) , but the irreducible polynomial \(f\) cannot lie in \(P[X]^p\) .

We thus have \(f' \neq 0\) and since the degree of \(f'\) is strictly less than that of f, we have \(f'(x) \neq 0\) . Now from \(f(x) = 0\) we deduce \(f'(x) \cdot dx = 0\) in \(\Omega_{\mathrm{P}}(\mathrm{B})\) , whence dx = 0, as we had to show.

* B) Suppose that A is etale; by the equivalence of a) and b) the algebra A is reduced, so there exist extensions \(L_{1}, \ldots, L_{n}\) of K such that A is isomorphic to \(L_{1} \times \cdots \times L_{n}\) (Prop. 5). Since A is etale, each of the extensions \(L_{i}\) is an etale algebra (V, p.~31, Cor.), and hence by definition is a separable algebraic extension of K.

The implication \(d) \Rightarrow a)\) follows from V, p.~31, Cor. *

COROLLARY. --- Suppose that K has characteristic \(p \neq 0\) . For A to be etale it is necessary and sufficient that \(A = K[A^{p}]\) . For every basis \((a_{i})_{i \in I}\) of A over K, the family \((a_{i}^{p})_{i \in I}\) is then a basis of A over K.

Let us choose an algebraic closure \(\Omega\) of K. Given two K-homomorphisms \(u\) and \(v\) of A into \(\Omega\) , if \(u\) and \(v\) have the same restriction to \(\mathbf{K}[\mathbf{A}^p]\) , we have

\[
u (x) ^ {p} = u \left(x ^ {p}\right) = v \left(x ^ {p}\right) = v (x) ^ {p}, \quad \text { whence } \quad u (x) = v (x)
\]

for all \(x \in \mathbf{A}\) . We thus have the inequality \([\mathbf{A} : \mathbf{K}]_s \leqslant [\mathbf{K}[\mathbf{A}^p] : \mathbf{K}]_s\) ; if \(\mathbf{A}\) is etale, we thus have

\[
[ \mathrm{A}: \mathrm{K} ] = [ \mathrm{A}: \mathrm{K} ] _ {s} \leqslant [ \mathrm{K} [ \mathrm{A} ^ {p} ]: \mathrm{K} ] _ {s} \leqslant [ \mathrm{K} [ \mathrm{A} ^ {p} ]: \mathrm{K} ],
\]

whence \(\mathbf{A} = \mathbf{K}[\mathbf{A}^p]\) .

Conversely suppose that we have \(\mathbf{A} = \mathbf{K}[\mathbf{A}^p]\) ; let \((a_i)_{i\in I}\) be a basis of \(\mathbf{A}\) over \(\mathbf{K}\) . By \(\mathbf{V}\) , p.~4, Prop. 2, \(b\) ), the family \((a_i^p)_{i\in I}\) generates the vector \(\mathbf{K}\) -space \(\mathbf{K}[\mathbf{A}^p]\) and since \(\mathbf{A} = \mathbf{K}[\mathbf{A}^p]\) is of finite dimension equal to the cardinal of \(\mathbf{I}\) , the family \((a_i^p)_{i\in I}\) is a basis of \(\mathbf{A}\) over \(\mathbf{K}\) . Let \(u\) be an element of \(\Omega \otimes_{\mathbf{K}} \mathbf{A}\) such that \(u^2 = 0\) , whence \(u^p = 0\) ; since \((a_i)_{i\in I}\) is a basis of \(\mathbf{A}\) over \(\mathbf{K}\) , there exists a family \((\lambda_i)_{i\in I}\) of elements of \(\Omega\) such that \(u = \sum_{i\in I} \lambda_i \otimes a_i\) , whence \(u^p = \sum_{i\in I} \lambda_i^p \otimes a_i^p\) . Since \((a_i^p)_{i\in I}\) is a basis of \(\mathbf{A}\) over \(\mathbf{K}\) and \(u^p = 0\) , we must have \(\lambda_i^p = 0\) , whence \(\lambda_i = 0\) for all \(i\in I\) , and so \(u = 0\) . This shows the ring \(\Omega \otimes_{\mathbf{K}} \mathbf{A}\) to be reduced; since \(\Omega\) is perfect, the algebra \(\mathbf{A}\) over \(\mathbf{K}\) is etale by Th. 4.

For another characterization of etale algebras see V, p.~48, Prop. 1.

